% Part: counterfactuals
% Chapter: minimal-change-semantics
% Section: antecedent-strengthening

\documentclass[../../../include/open-logic-section]{subfiles}

\begin{document}

\olfileid{cnt}{min}{agg}

\olsection{পূর্বপক্ষ শক্তিশালীকরণ}

% BN-SRC-425: মহাশূন্যের উদাহরণে source-এর জ্বালানোর বদলে ঘষার পূর্বপক্ষ চাই।
``পূর্বপক্ষকে শক্তিশালী করা'' বলতে
$!A \lif !C \Entails (!A \land !B) \lif !C$ অনুমিতিটি বোঝায়।
বস্তুগত শর্তবচনের ক্ষেত্রে এটি বৈধ, কিন্তু প্রতিবাস্তব
শর্তবচনের ক্ষেত্রে অবৈধ। ধরা যাক, এই দেশলাই-কাঠিটি
ঘষলে জ্বলে উঠত—এটি সত্য। (অর্থাৎ কাঠিটিতে বা বাক্সের
ঘষার পৃষ্ঠে কোনও খুঁত নেই, আমি কাঠিটি ভেঙে ফেলব না
ইত্যাদি।) কিন্তু মহাশূন্যে এই কাঠিটি ঘষলে জ্বলে উঠত—এটি
সত্য নয়। তাই নিচের অনুমিতি অবৈধ:
\begin{quote}
  দেশলাই-কাঠিটি ঘষা হলে জ্বলে উঠত।

  অতএব, মহাশূন্যে দেশলাই-কাঠিটি ঘষা হলে জ্বলে উঠত।
\end{quote}

শর্তবচন সম্পর্কে লুইস--স্টালনেকারের ব্যাখ্যা এর কারণ দেখায়:
মহাশূন্যে আমি দেশলাই-কাঠিটি ঘষি এমন নিকটতম জগৎটি,
কেবল কাঠিটি ঘষি এমন নিকটতম জগৎটির তুলনায় বাস্তব জগৎ
থেকে অনেক দূরের। তাই দ্বিতীয় জগতে কাঠিটি জ্বললেও
প্রথম জগতে জ্বলে না। স্বজ্ঞাগত প্রত্যাশাও তা-ই।

\begin{ex}
  গোলক অর্থতত্ত্বে অনুমিতিটি অবৈধ, অর্থাৎ $p
  \cif r \Entails/ (p \land q) \cif r$। ধরা যাক,
  $\mModel{M} = \tuple{W, O, V}$ মডেলে
  $W = \{w, w_1, w_2\}$,
  $O_w = \{\{w\}, \{w, w_1\}, \{w, w_1, w_2\}\}$,
  $V(p) = \{w_1, w_2\}$, $V(q) = \{w_2\}$
  এবং $V(r) = \{w_1\}$।
  % BN-SRC-773: অন্য জগতগুলিতেও O মোট অপেক্ষক করতে গোলকব্যবস্থা নির্দিষ্ট।
  অন্য দুই জগতে প্রতিটির নিজের এককবিশিষ্ট সেটমাত্র দিয়ে
  গোলকব্যবস্থা নেওয়া যায়; এতে মূল জগতের নিচের মূল্যায়ন বদলায় না।
  $p$-স্বীকারক গোলক
  $S = \{w, w_1\}$-এর প্রতিটি জগতে $p \lif r$ সত্য;
  তাই $\mSat{M}{p \cif r}[w]$। আবার $(p \land q)$-স্বীকারক
  গোলক $S' = \{w, w_1, w_2\}$ আছে, কিন্তু
  $\mSat/{M}{(p \land q) \lif r}[w_2]$। এখানে এই
  পূর্বপক্ষ-স্বীকারক গোলকটিই একমাত্র; ছোট গোলকগুলিতে পূর্বপক্ষ
  সত্য কোনো জগৎ নেই। তাই
  $\mSat/{M}{(p \land q) \cif r}[w]$
  (দেখুন \olref{fig:antecedent-strengthening})।

\begin{figure}
\begin{center}
\begin{tikzpicture}[layerwidth=1.5,scale=.6]\tiny
  \spheresystem{3}
  \propositionintersect{-10}{3}{30}{5}
  \begin{scope}
    \clip plot[smooth,tension=1] coordinates
          {(0,4.5) (30:.95) (4.5,-3.5)} -- (4.5,4.5) ;
    \spherefill{2}
  \end{scope}
  \draw[proposition] plot[dashed,smooth,tension=1] coordinates
          {(0,4.5) (30:.95) (4.5,-3.5)} ;
  \proposition{30}{2}{30}{5}
  \path (0,0) node[world] {$w$};
  \spherepos{30}{2}{node[world] {$w_1$}}
  \spherepos{-10}{3}{node[world] {$w_2$}}
  \spherepos{-10}{4}{node {$q$}}
  \spherepos{30}{4}{node {$r$}}
  \draw (1,4.5) node[below] {$p$};
\end{tikzpicture}
\caption{পূর্বপক্ষ শক্তিশালীকরণের প্রত্যুদাহরণ}
\ollabel{fig:antecedent-strengthening}
\end{center}
\end{figure}
\end{ex}

\end{document}
